% § 2 — Background and Related Work
\section{Background and Related Work}
\label{sec:related}
\textbf{Unfair-term detection.} \corpus{}~\cite{lippi2019claudette,lagioia2019deep} labels sentences of 50 ToS with eight categories---arbitration, unilateral change, content removal, jurisdiction, choice of law, limitation of liability, unilateral termination and contract by using---whose definitions are themselves derived from the Annex and from doctrine~\cite{micklitz2017empire}. The corpus has been extended to several languages~\cite{drawzeski2021corpus,galassi2024multilingual}, included in LexGLUE~\cite{chalkidis2022lexglue}, and used to compare fine-tuned encoders with prompted LLMs~\cite{panarelli2025worth}. Related resources annotate clause validity in German consumer contracts~\cite{braun2024agbde} and potentially abusive clauses in Chilean ToS~\cite{loffler2025chilean}. All of them predict; none of them cites the legal ground of a prediction.

\textbf{Legal knowledge representation.} Rule languages such as LegalRuleML~\cite{athan2015legalruleml} and contract specification languages such as Symboleo~\cite{sharifi2020symboleo} formalise norms as deontic statements with parties, actions and conditions; they are designed for authoring and reasoning, not for being populated from natural-language contracts at scale. Deontic modality extraction~\cite{chalkidis2018obligation,sancheti2022lexdemod} recovers obligations, permissions and prohibitions from text, with macro-F1 in the 0.6--0.8 range depending on the setting---good enough to be an input signal, not good enough to be a ground truth. Large language models annotate legal text reliably enough in zero-shot settings to serve as extractors~\cite{savelka2023unreasonable}; we use one under a closed schema, as an extractor, never as a judge of unfairness. Our templates sit between these two traditions: a fixed, shallow structure that lawyers can validate clause by clause and that queries can be written against.

\textbf{Queries over contracts.} ContractNLI~\cite{koreeda2021contractnli} evaluates document-level hypotheses (``the receiving party may share confidential information with employees'') by learned entailment; CUAD~\cite{hendrycks2021cuad} retrieves clause types for due diligence. Both pose the question we pose---does this contract contain a provision of this kind?---but answer it with a learned model. We answer it by executing an explicit condition and returning the template fields that satisfy it.

\textbf{Explanation in AI and law.} The field has long argued that legal decision support must explain in legal terms~\cite{atkinson2020explanation}. For unfair-term detection, explanations have so far been rationales selected by attention~\cite{ruggeri2022memory} or generated by a model~\cite{liepina2020claudettetool}. A query that fires because a termination clause grants the provider a discretionary right without notice explains itself by construction, and the explanation can be wrong only in a way a lawyer can check.
